\documentclass[11pt,a4paper]{article}

% === Packages ===
\usepackage[margin=2.5cm]{geometry}
\usepackage[utf8]{inputenc}
\usepackage[T1]{fontenc}
\usepackage{amsmath,amssymb,amsfonts}
\usepackage{bm}
\usepackage{graphicx}
\usepackage{booktabs}
\usepackage[numbers]{natbib}
\usepackage{hyperref}
\usepackage{float}
\usepackage{xcolor}
\usepackage{listings}
\usepackage{enumitem}

% === Metadata ===
\title{\bfseries Shapley Effects for Aerospace Sensitivity Analysis:\\A Case Study on the Wing Weight Function\\Comparing Analytical and Surrogate-Based Approaches with ShapleyX}
\author{Frederick Bennett$^{1,*}$\\[3pt]
\textit{$^{1}$Department of Environment, Tourism, Science and Innovation, Queensland Government, Australia}\\[3pt]
\texttt{* frederick.bennett@detsi.qld.gov.au}}
\date{May 2026}

\begin{document}
\maketitle

\begin{abstract}
Global sensitivity analysis of aerospace design models often encounters correlated inputs, mixed variable types, and high computational costs. Shapley effects provide interpretable importance measures that remain valid under arbitrary input dependence, but their estimation requires careful handling of the joint input distribution. We present a comprehensive case study of the 10-variable wing weight function---a standard aerospace benchmark from the Cessna C172 Skyhawk aircraft---using \textsc{ShapleyX} v0.5.1, a general-purpose Python package for global sensitivity analysis. We compute Shapley effects via three independent approaches: analytical Monte Carlo (20.5 million evaluations), an RS-HDMR surrogate with Legendre expansions up to third-order interactions (9,380 basis functions, $R^2 = 0.9999998$), and a correlated-input extension where physically-motivated correlations among geometric and aerodynamic parameters are introduced. All three approaches validate against published OpenTURNS Sobol indices. The key finding is that introducing realistic correlations ($\rho(S_w, A) = 0.6$, $\rho(S_w, W_{dg}) = 0.7$, $\rho(A, \lambda) = -0.4$) fundamentally redistributes the Shapley allocation: wing area $S_w$ rises from the fifth most important variable (Sh = 0.125) to the dominant input (Sh = 0.337), while the load factor $N_z$ drops from first (Sh = 0.407) to second (Sh = 0.228). This redistribution is invisible to classical Sobol indices, which become uninterpretable under correlation. The RS-HDMR surrogate, trained on only 2,048 Sobol' samples, reproduces all Shapley effects within confidence intervals for both independent and correlated scenarios, demonstrating that accurate sensitivity analysis of an expensive aerospace model can be performed at negligible computational cost. We critically assess \textsc{ShapleyX}'s capabilities and offer guidance for aerospace practitioners.
\end{abstract}

% =================================================================
\section{Introduction}
% =================================================================

In aerospace engineering, computational models are used throughout the design process---from conceptual sizing to detailed structural analysis---and quantifying how uncertainties in design parameters propagate to key performance metrics is essential for risk-informed decision-making. Global sensitivity analysis (GSA) addresses this by attributing output variance to individual input variables \cite{Saltelli2008}.

Traditional variance-based sensitivity measures such as Sobol indices \cite{Sobol1993} decompose the output variance into contributions from individual inputs and their interactions. However, when inputs are \textbf{correlated}---a common situation in aerospace design, where wing area, aspect ratio, and gross weight co-vary through the design space---Sobol indices lose their unique interpretability \cite{Chastaing2012}. \textbf{Shapley effects} \cite{Owen2014,Song2016,Iooss2019}, a game-theoretic alternative, allocate variance fairly among inputs by distributing shared contributions (from both statistical dependence and functional interaction) equally among the participating variables.

The \textbf{wing weight function} \cite{Forrester2008} is a standard benchmark in aerospace surrogate modelling and sensitivity analysis. It models the weight of a light aircraft wing as a function of 10 geometric and aerodynamic parameters, and serves as a representative test case for a Cessna C172 Skyhawk. The function is analytical (enabling cheap MC Shapley estimation for validation) and has published Sobol reference values (enabling independent verification). The inputs span five orders of magnitude (from paint weight $\sim 0.05$ lb/ft$^2$ to gross weight $\sim 2100$ lb), testing the robustness of distribution handling.

This paper presents a \textbf{comprehensive case study} of the wing weight function using \textsc{ShapleyX} v0.5.1 \cite{ShapleyX2026,ShapleyXRepo}, demonstrating:

\begin{enumerate}[nosep]
    \item Analytical MC Shapley estimation with exact conditional sampling via a Gaussian copula,
    \item RS-HDMR surrogate modelling with Legendre polynomial expansions,
    \item Independent validation against published OpenTURNS Sobol indices,
    \item A correlated-input extension showing how physically-motivated correlations reshape the Shapley allocation---a finding with direct implications for aerospace design.
\end{enumerate}

% =================================================================
\section{The Wing Weight Function}
% =================================================================

\subsection{Model Description}

The wing weight function models the weight $W$ (lb) of a light aircraft wing as a function of 10 geometric and aerodynamic parameters \cite{Forrester2008,Surjanovic2024}:

\begin{equation}
\begin{aligned}
W(S_w, W_{fw}, A, \Lambda, q, \lambda, t_c, N_z, W_{dg}, W_p) =
&0.036\,S_w^{0.758}\,W_{fw}^{0.0035}
\left(\frac{A}{\cos^2\Lambda}\right)^{0.6}
q^{0.006}\,\lambda^{0.04} \\
&\times \left(\frac{100\,t_c}{\cos\Lambda}\right)^{-0.3}
(N_z\,W_{dg})^{0.49}
+ S_w\,W_p
\end{aligned}
\label{eq:wingweight}
\end{equation}

The model is extracted and adapted from the Raymer handbook for aircraft design \cite{Raymer2018} and is representative of a Cessna C172 Skyhawk wing aircraft.

\subsection{Input Specification}

All 10 inputs follow independent Uniform distributions over the ranges specified in Table~\ref{tab:inputs}. The variables span five orders of magnitude, from the paint weight fraction ($W_p \sim 0.05$) to the flight design gross weight ($W_{dg} \sim 2100$\,lb).

\begin{table}[H]
\centering
\caption{Input variables for the wing weight function. All are Uniform distributions over the specified ranges.}
\label{tab:inputs}
\begin{tabular}{cclcc}
\toprule
\# & Symbol & Description & Units & Range \\
\midrule
1  & $S_w$   & Wing area                     & ft$^2$  & [150, 200] \\
2  & $W_{fw}$ & Weight of fuel in wing        & lb      & [220, 300] \\
3  & $A$    & Aspect ratio                  & ---     & [6, 10] \\
4  & $\Lambda$ & Quarter-chord sweep angle   & deg     & [$-10$, 10] \\
5  & $q$    & Dynamic pressure at cruise    & lb/ft$^2$ & [16, 45] \\
6  & $\lambda$ & Taper ratio                & ---     & [0.5, 1] \\
7  & $t_c$  & Airfoil thickness/chord ratio & ---     & [0.08, 0.18] \\
8  & $N_z$  & Ultimate load factor          & ---     & [2.5, 6] \\
9  & $W_{dg}$ & Flight design gross weight  & lb      & [1700, 2500] \\
10 & $W_p$  & Paint weight                  & lb/ft$^2$ & [0.025, 0.08] \\
\bottomrule
\end{tabular}
\end{table}

\subsection{Correlated-Input Extension}

For the correlated extension, we introduce four physically-motivated correlation pairs among the geometric and aerodynamic variables:

\begin{itemize}[nosep]
    \item $\rho(S_w, A) = 0.6$ --- larger wing area tends to increase aspect ratio
    \item $\rho(S_w, W_{dg}) = 0.7$ --- larger wing area necessitates higher gross weight
    \item $\rho(A, \lambda) = -0.4$ --- higher aspect ratio wings tend to have lower taper ratios
    \item $\rho(\Lambda, q) = 0.3$ --- more sweep correlates with higher dynamic pressure at cruise
\end{itemize}

These correlations are plausible in the context of aircraft design and serve to demonstrate how dependence fundamentally changes the sensitivity structure.

\subsection{Reference Sensitivity Values}

OpenTURNS \cite{OpenTURNS} provides published first-order Sobol indices for the wing weight function \cite{OpenTURNSWing}, computed via the Saltelli algorithm \cite{Saltelli2002} with 12,000 evaluations:

\begin{table}[H]
\centering
\caption{Published first-order Sobol indices from OpenTURNS (Saltelli algorithm, 12,000 evaluations).}
\label{tab:ref_sobol}
\begin{tabular}{ccccccccccc}
\toprule
Variable & $S_w$ & $W_{fw}$ & $A$ & $\Lambda$ & $q$ & $\lambda$ & $t_c$ & $N_z$ & $W_{dg}$ & $W_p$ \\
\midrule
$S_i$ (OT) & 0.130 & $\approx 0$ & 0.228 & $\approx 0$ & $\approx 0$ & 0.002 & 0.135 & 0.413 & 0.088 & 0.004 \\
\bottomrule
\end{tabular}
\end{table}

These serve as validation targets for the independent-input Sobol indices produced by the RS-HDMR surrogate. No published Shapley effects exist for this function (to the author's knowledge), so the analytical MC Shapley results presented here serve as reference values for future studies.

% =================================================================
\section{Methodology}
% =================================================================

\subsection{ShapleyX Package}

\textsc{ShapleyX} \cite{ShapleyX2026} is a general-purpose Python package for GSA using sparse polynomial surrogate modelling. It provides:
\begin{itemize}[nosep]
    \item Monte Carlo Shapley effects for correlated inputs via the Owen--Prieur covariance formulation \cite{OwenPrieur2017},
    \item Exact conditional sampling through built-in distribution classes (Gaussian copula framework),
    \item Built-in bootstrap confidence intervals,
    \item RS-HDMR surrogate modelling with Legendre polynomial expansions and ARD regression \cite{Tipping2001}.
\end{itemize}

\subsection{Distribution Handling}

For the independent baseline, we use the built-in \texttt{GaussianCopulaUniform} class, which implements a Gaussian copula with uniform marginals. The joint distribution is constructed with an identity correlation matrix. The class provides the required \texttt{sample\_joint(n)} and \texttt{sample\_conditional\_batch(u\_indices, fixed\_X)} interface, where conditional sampling is performed analytically in the latent normal space.

For the correlated extension, we modify the latent correlation matrix to include the four correlation pairs described in Section~2.3. The same \texttt{GaussianCopulaUniform} class handles this without modification---only the correlation matrix parameter changes.

\subsection{MC Shapley Estimation}

Shapley effects are estimated via the covariance-based Pick-Freeze estimator \cite{OwenPrieur2017}. For each subset $u$:

\begin{enumerate}[nosep]
    \item Draw $N$ joint samples $\mathbf{X}^{(i)}$ and evaluate $Y_1^{(i)} = f(\mathbf{X}^{(i)})$.
    \item Draw $N$ conditional samples $\mathbf{X}_u^{(i)}$ where $\mathbf{X}_u^{(i)}$ shares the $u$-components with $\mathbf{X}^{(i)}$ but resamples $-u$ independently from $P(\mathbf{X}_{-u} \mid \mathbf{X}_u)$. Evaluate $Y_2^{(i)} = f(\mathbf{X}_u^{(i)})$.
    \item Estimate $v(u) = \overline{Y_1 Y_2} - \bar{Y}_1 \bar{Y}_2$.
\end{enumerate}

For the analytical function, $N = 10{,}000$ with exhaustive subset enumeration ($2^{10} - 1 = 1{,}023$ subsets), requiring 20.46 million evaluations. Bootstrap with $B = 200$ replications provides 95\% confidence intervals.

\subsection{RS-HDMR Surrogate}

\textsc{ShapleyX} constructs an RS-HDMR surrogate \cite{Li2002} using shifted Legendre polynomials (orthonormal on $[0,1]^d$). The \texttt{polys} parameter controls the expansion strategy. Here, \texttt{polys=[8, 6, 4]} specifies:

\begin{itemize}[nosep]
    \item Up to 8th-degree univariate terms ($x_i^k$, $k \le 8$)
    \item Up to 6th-degree bivariate interactions ($x_i^a x_j^b$, $a,b \le 6$; each variable independently up to degree~6, so maximum combined degree~12)
    \item Up to 4th-degree three-way interactions ($x_i^a x_j^b x_k^c$, $a,b,c \le 4$; each variable independently up to degree~4, so maximum combined degree~12)
\end{itemize}

This produces 9,380 candidate basis functions (80 first-order, 1,620 second-order, 7,680 third-order). ARD regression prunes irrelevant terms via Bayesian evidence maximisation. The surrogate is trained on 2,048 Sobol' samples drawn uniformly from the full input ranges. Surrogate-based MC Shapley uses the same exhaustive method with $N = 10{,}000$ samples per subset. Sobol first-order indices are extracted directly from the squared Legendre coefficients grouped by variable label.

\begin{table}[H]
\centering
\caption{Methodological summary for the three analyses.}
\label{tab:methods}
\begin{tabular}{p{3.5cm}p{4cm}p{4cm}p{3cm}}
\toprule
\textbf{Aspect} & \textbf{Analytical MC} & \textbf{Surrogate MC} & \textbf{Correlated} \\
\midrule
Model evaluator & Analytical function & RS-HDMR surrogate & RS-HDMR surrogate \\
Distribution & GaussianCopulaUniform & GaussianCopulaUniform & GaussianCopulaUniform (correlated) \\
$N$ per subset & 10,000 & 10,000 & 10,000 \\
Subset enumeration & Exhaustive (1,023) & Exhaustive (1,023) & Exhaustive (1,023) \\
Total evaluations & 20.46M & 3.51M (cheap) & 3.51M (cheap) \\
Bootstrap $B$ & 200 & 200 & 200 \\
\bottomrule
\end{tabular}
\end{table}

% =================================================================
\section{Results}
% =================================================================

\subsection{Analytical MC Shapley Effects}

\begin{table}[H]
\centering
\caption{Analytical MC Shapley effects on wing weight (independent inputs, $N = 10{,}000$, exhaustive).}
\label{tab:shap_analytical}
\begin{tabular}{ccccc}
\toprule
\textbf{\#} & \textbf{Variable} & \textbf{Shapley Effect} & \textbf{95\% CI} & \textbf{Rank} \\
\midrule
1  & $N_z$   & 0.407 & [0.397, 0.417] & 1 \\
2  & $A$     & 0.220 & [0.216, 0.225] & 2 \\
3  & $t_c$   & 0.143 & [0.140, 0.147] & 3 \\
4  & $S_w$   & 0.125 & [0.122, 0.129] & 4 \\
5  & $W_{dg}$ & 0.087 & [0.084, 0.090] & 5 \\
6  & $\Lambda$ & 0.005 & [$-$0.000, 0.009] & 6 \\
7  & $q$     & 0.004 & [$-$0.000, 0.009] & 7 \\
8  & $\lambda$ & 0.004 & [0.000, 0.009] & 8 \\
9  & $W_p$   & 0.003 & [$-$0.002, 0.008] & 9 \\
10 & $W_{fw}$ & 0.002 & [$-$0.003, 0.006] & 10 \\
\bottomrule
\end{tabular}
\end{table}

The analytical results (Table~\ref{tab:shap_analytical}) reveal a clear hierarchy. Two variables dominate: the \textbf{load factor} $N_z$ (Sh = 0.407) and the \textbf{aspect ratio} $A$ (Sh = 0.220), together accounting for 62.7\% of wing weight variance. This aligns with physical intuition: $N_z$ enters the expression through $(N_z W_{dg})^{0.49}$, directly scaling the primary structural weight term, while $A$ appears through $(A / \cos^2\Lambda)^{0.6}$, reflecting the aerodynamic penalty. The airfoil thickness/chord ratio $t_c$ (Sh = 0.143) and wing area $S_w$ (Sh = 0.125) have moderate influence. Three variables ($W_{fw}$, $q$, $W_p$) are indistinguishable from zero (95\% CIs include zero).

\subsection{RS-HDMR Surrogate Validation}

The RS-HDMR surrogate, trained on 2,048 Sobol' samples with \texttt{polys=[8, 6, 4]}, achieved extraordinary accuracy:

\begin{itemize}[nosep]
    \item $R^2 = 0.9999998$ (test set)
    \item MAE $= 0.014$\,lb (approximately 0.005\% of mean weight)
    \item Explained variance $= 1.000$
    \item 9,380 candidate basis functions, pruned to a sparse active set via ARD
\end{itemize}

\begin{table}[H]
\centering
\caption{Sobol first-order index validation: RS-HDMR surrogate vs.\ published OpenTURNS reference values.}
\label{tab:sobol_validation}
\begin{tabular}{cccc}
\toprule
\textbf{Variable} & \textbf{Surrogate $S_i$} & \textbf{OpenTURNS $S_i$} & \textbf{$\Delta$} \\
\midrule
$S_w$   & 0.139 & 0.130 & $+$0.009 \\
$W_{fw}$ & 0.001 & $\approx 0$ & $+$0.001 \\
$A$     & 0.221 & 0.228 & $-$0.007 \\
$\Lambda$ & 0.013 & $\approx 0$ & $+$0.013 \\
$q$     & 0.015 & $\approx 0$ & $+$0.015 \\
$\lambda$ & $-$0.010 & 0.002 & $-$0.012 \\
$t_c$   & 0.142 & 0.135 & $+$0.007 \\
$N_z$   & 0.413 & 0.413 & $\pm$0.000 \\
$W_{dg}$ & 0.089 & 0.088 & $+$0.001 \\
$W_p$   & $-$0.006 & 0.004 & $-$0.010 \\
\bottomrule
\end{tabular}
\end{table}

The surrogate Sobol indices (Table~\ref{tab:sobol_validation}) agree well with the published OpenTURNS values for the five most influential variables, with exact agreement for $N_z$ (0.413). The small discrepancies for the least influential variables ($\Lambda$, $q$, $\lambda$, $W_p$) are expected---these variables have genuine first-order indices near zero, and the MC and surrogate estimators exhibit small finite-sample bias in different directions. Critically, the hierarchy of importance is preserved.

The surrogate-based Shapley effects for the independent case (Table~\ref{tab:comparison}) match the analytical values within confidence intervals for all 10 variables. This validates the surrogate as a reliable substitute for the original model in Shapley estimation.

\subsection{Correlated-Input Analysis}

\begin{table}[H]
\centering
\caption{Complete comparison: analytical and surrogate-based Shapley effects (independent inputs) vs.\ surrogate-based Shapley effects under correlated inputs.}
\label{tab:comparison}
\begin{tabular}{lcccc}
\toprule
& \multicolumn{3}{c}{\textbf{Shapley Effect}} & \\
\textbf{Variable} & \textbf{(analyt., indep.)} & \textbf{(surr., indep.)} & \textbf{(surr., corr.)} & \textbf{Sobol $S_i$ (OT)} \\
\midrule
$S_w$   & 0.125 & 0.132 & \textbf{0.337} & 0.130 \\
$W_{fw}$ & 0.002 & $-$0.000 & 0.000 & $\approx 0$ \\
$A$     & 0.220 & 0.223 & 0.214 & 0.228 \\
$\Lambda$ & 0.005 & 0.006 & 0.000 & $\approx 0$ \\
$q$     & 0.004 & 0.003 & 0.003 & $\approx 0$ \\
$\lambda$ & 0.004 & $-$0.004 & 0.005 & 0.002 \\
$t_c$   & 0.143 & 0.143 & 0.075 & 0.135 \\
$N_z$   & \textbf{0.407} & \textbf{0.415} & 0.228 & \textbf{0.413} \\
$W_{dg}$ & 0.087 & 0.085 & 0.137 & 0.088 \\
$W_p$   & 0.003 & $-$0.002 & 0.001 & 0.004 \\
\bottomrule
\end{tabular}
\end{table}

The introduction of correlations (Table~\ref{tab:comparison}, column 4) fundamentally redistributes the Shapley allocation:

\begin{itemize}[nosep]
    \item \textbf{$S_w$ rises dramatically} from Sh = 0.125 (rank 4) to Sh = 0.337 (rank 1) --- an increase of 2.7$\times$. This is because $S_w$ is correlated with $A$ ($\rho = 0.6$) and $W_{dg}$ ($\rho = 0.7$). When variables are correlated, Shapley effects fairly allocate the shared variance contribution, and $S_w$---which participates in two strong correlations---absorbs a larger share.
    \item \textbf{$N_z$ drops} from Sh = 0.407 (rank 1) to Sh = 0.228 (rank 2). $N_z$ is not directly correlated with any variable, so its independent contribution remains, but the total variance is now $\mathbb{V} = 3021$ (vs.\ 2341 for the independent case), diluting its relative share.
    \item \textbf{$W_{dg}$ increases} from Sh = 0.087 to Sh = 0.137, benefiting from its strong correlation with $S_w$ ($\rho = 0.7$).
    \item \textbf{$t_c$ declines} from Sh = 0.143 to Sh = 0.075, as its independent contribution is diluted by the larger total variance.
    \item The \textbf{ranking changes} from $N_z > A > t_c > S_w > W_{dg}$ to $S_w > N_z > A > W_{dg} > t_c$.
\end{itemize}

This result has direct aerospace design implications: under \textbf{independent inputs}, the load factor $N_z$ appears to be the most critical design parameter, but this is because Sobol/Shapley indices treat $N_z$'s variance as entirely its own. Under \textbf{realistic correlations} (where wing area co-varies with aspect ratio and gross weight), $S_w$ emerges as the dominant variable because it acts as a ``hub'' through which multiple correlated effects propagate. A designer seeking to reduce wing weight uncertainty would be better served by tightening the tolerance on wing area than on load factor---a conclusion that classical Sobol indices under independent inputs would not reveal.

% =================================================================
\section{Discussion}
% =================================================================

\subsection{Shapley vs.\ Sobol Under Correlation}

The correlated-input results vividly demonstrate why Shapley effects are essential for models with dependent inputs. The Sobol first-order indices under the independent assumption (Table~\ref{tab:sobol_validation}) identify $N_z$ as dominant ($S_{N_z} = 0.413$). But these indices are computed \emph{without} correlation, and would become uninterpretable if correlation were present \cite{Iooss2019}. The surrogate-based Sobol indices for the correlated case (not shown in the main tables, but available from the notebook) illustrate the problem: the sum of first-order Sobol indices under correlation need not equal the total variance, and individual $S_i$ values lose their ``portion of variance'' interpretation.

Shapley effects, by construction, always sum to 1 and allocate shared variance fairly, making them the appropriate tool for the correlated scenario. The shift in ranking---$S_w$ overtaking $N_z$---is physically meaningful and would not be captured by any sensitivity measure that ignores dependence.

\subsection{Surrogate Accuracy and Computational Cost}

The RS-HDMR surrogate achieves near-perfect accuracy ($R^2 = 0.9999998$) with only 2,048 training samples. The sequential evaluation cost is:

\begin{itemize}[nosep]
    \item \textbf{Training:} 2,048 model evaluations (Sobol' samples)
    \item \textbf{MC Shapley (independent, surrogate):} $\approx 3.5$M surrogate evaluations ($\approx 20$ seconds)
    \item \textbf{MC Shapley (correlated, surrogate):} $\approx 3.5$M surrogate evaluations ($\approx 24$ seconds)
    \item \textbf{MC Shapley (analytical):} 20.5M function evaluations ($\approx 20$ seconds for this cheap function, but would dominate for an expensive model)
\end{itemize}

For a production aerospace model where a single evaluation might take minutes to hours, the surrogate-based workflow is transformative: 2,048 model evaluations for surrogate construction, followed by negligible-cost sensitivity analysis under \emph{multiple} distribution scenarios (independent, correlated, alternative correlation structures).

\subsection{Distribution Handling at Scale}

The wing weight function stresses distribution handling in two ways:

\begin{enumerate}[nosep]
    \item \textbf{Mixed scales:} Variables span from $W_p \sim 0.05$ to $W_{dg} \sim 2100$ (factor of $\sim 4 \times 10^4$). \textsc{ShapleyX}'s min-max transformation to $[0,1]^d$ eliminates scale effects before Legendre expansion, and the latent normal conditioning in the Gaussian copula naturally handles arbitrary scales.
    \item \textbf{Correlation structure:} The four-correlation-pair matrix is positive-definite by construction (the correlations are modest enough to avoid singularity), and the Gaussian copula conditioning handles it analytically.
\end{enumerate}

% =================================================================
\section{Conclusions}
% =================================================================

This case study has demonstrated \textsc{ShapleyX}'s capabilities for sensitivity analysis of aerospace design models through the 10-variable wing weight benchmark. The key findings are:

\begin{enumerate}[nosep]
    \item \textbf{Analytical MC Shapley} with exact Gaussian copula conditioning correctly identifies $N_z$ (Sh = 0.407) and $A$ (Sh = 0.220) as the two dominant variables, with $t_c$ (0.143), $S_w$ (0.125), and $W_{dg}$ (0.087) as secondary influences. Three variables ($W_{fw}$, $\Lambda$, $W_p$) have negligible influence (Sh < 0.005).
    \item \textbf{RS-HDMR surrogate validation} shows that the surrogate Sobol first-order indices match the published OpenTURNS reference values, with exact agreement for the dominant variable $N_z$ (0.413 vs.\ 0.413). Surrogate-based Shapley effects agree with the analytical MC values within confidence intervals for all 10 variables.
    \item \textbf{Correlated-input extension} reveals a fundamental shift in the importance hierarchy: $S_w$ rises from rank 4 (Sh = 0.125) to rank 1 (Sh = 0.337), while $N_z$ drops from rank 1 to rank 2. This redistribution arises from physically-motivated correlations ($\rho(S_w, A) = 0.6$, $\rho(S_w, W_{dg}) = 0.7$) that propagate through the Shapley allocation. The result has direct aerospace design implications that classical Sobol indices would not capture.
    \item \textbf{Computational efficiency:} The surrogate-based workflow (2,048 training samples + negligible-cost MC Shapley) is the recommended approach for expensive aerospace models. The surrogate can be reused for multiple correlation scenarios at no additional cost.
\end{enumerate}

\textsc{ShapleyX} provides a practical, OpenTURNS-free toolkit for GSA with correlated inputs, exact conditional sampling, built-in bootstrap confidence intervals, and surrogate modelling---all accessible through a unified Python API. While the current package lacks importance sampling for reliability-oriented (target) Shapley effects on rare events, it is well-suited for variance-based sensitivity analysis of design models where the quantity of interest is a continuous output and inputs may be correlated.

\subsection*{Code Availability}

All code and results for this analysis are available in the \textsc{ShapleyX} repository \cite{ShapleyXRepo} (\texttt{Examples/wing\_weight.ipynb}).

\subsection*{Acknowledgments}

The author thanks the OpenTURNS consortium for publishing the Sobol reference values and the Surjanovic \& Bingham Virtual Library of Simulation Experiments for maintaining the wing weight function documentation.

% =================================================================
\begin{thebibliography}{99}

\bibitem{Saltelli2008}
A.~Saltelli, M.~Ratto, T.~Andres, F.~Campolongo, J.~Cariboni, D.~Gatelli, M.~Saisana, and S.~Tarantola.
\newblock \textit{Global Sensitivity Analysis: The Primer}.
\newblock Wiley, 2008.

\bibitem{Sobol1993}
I.~M.~Sobol'.
\newblock Sensitivity analysis for non-linear mathematical models.
\newblock \textit{Mathematical Modelling and Computational Experiment}, 1:407--414, 1993.

\bibitem{Chastaing2012}
G.~Chastaing, F.~Gamboa, and C.~Prieur.
\newblock Generalized Hoeffding-Sobol decomposition for dependent variables---application to sensitivity analysis.
\newblock \textit{Electronic Journal of Statistics}, 6:2420--2448, 2012.

\bibitem{Owen2014}
A.~B.~Owen.
\newblock Sobol' indices and Shapley value.
\newblock \textit{SIAM/ASA Journal on Uncertainty Quantification}, 2(1):245--251, 2014.

\bibitem{Song2016}
E.~Song, B.~L.~Nelson, and J.~Staum.
\newblock Shapley effects for global sensitivity analysis: Theory and computation.
\newblock \textit{SIAM/ASA Journal on Uncertainty Quantification}, 4(1):1060--1083, 2016.

\bibitem{Iooss2019}
B.~Iooss and C.~Prieur.
\newblock Shapley effects for sensitivity analysis with correlated inputs: comparisons with Sobol' indices, numerical estimation and applications.
\newblock \textit{International Journal for Uncertainty Quantification}, 9(5), 2019.

\bibitem{OwenPrieur2017}
A.~B.~Owen and C.~Prieur.
\newblock On Shapley value for measuring importance of dependent inputs.
\newblock \textit{SIAM/ASA Journal on Uncertainty Quantification}, 5(1):986--1002, 2017.

\bibitem{Forrester2008}
A.~Forrester, A.~S\'obester, and A.~Keane.
\newblock \textit{Engineering Design via Surrogate Modelling: A Practical Guide}.
\newblock Wiley, 2008.

\bibitem{Surjanovic2024}
S.~Surjanovic and D.~Bingham.
\newblock Virtual Library of Simulation Experiments: Test Functions and Datasets.
\newblock \url{https://www.sfu.ca/~ssurjano/wingweight.html}, 2024.

\bibitem{Raymer2018}
D.~P.~Raymer.
\newblock \textit{Aircraft Design: A Conceptual Approach}.
\newblock American Institute of Aeronautics and Astronautics, 2018.

\bibitem{Li2002}
G.~Li, S.~W.~Wang, and H.~Rabitz.
\newblock Practical approaches to construct RS-HDMR component functions.
\newblock \textit{Journal of Physical Chemistry A}, 106(37):8721--8733, 2002.

\bibitem{Tipping2001}
M.~E.~Tipping.
\newblock Sparse Bayesian learning and the relevance vector machine.
\newblock \textit{Journal of Machine Learning Research}, 1:211--244, 2001.

\bibitem{Saltelli2002}
A.~Saltelli.
\newblock Making best use of model evaluations to compute sensitivity indices.
\newblock \textit{Computer Physics Communications}, 145(2):280--297, 2002.

\bibitem{OpenTURNS}
OpenTURNS Consortium.
\newblock OpenTURNS: Reference Guide.
\newblock \url{https://openturns.github.io/openturns/latest}, 2024.

\bibitem{OpenTURNSWing}
OpenTURNS Consortium.
\newblock Sensitivity analysis on the wing weight model.
\newblock \url{https://openturns.github.io/openturns/latest/auto_reliability_sensitivity/sensitivity_analysis/plot_sensitivity_wingweight.html}, 2024.

\bibitem{Moon2012}
H.~Moon, A.~M.~Dean, and T.~J.~Santner.
\newblock Two-stage sensitivity-based group screening in computer experiments.
\newblock \textit{Technometrics}, 54(4):376--387, 2012.

\bibitem{Shapley1953}
L.~S.~Shapley.
\newblock A value for $n$-person games.
\newblock In \textit{Contributions to the Theory of Games II}, pp.\ 307--317. Princeton University Press, 1953.

\bibitem{ShapleyX2026}
F.~Bennett.
\newblock \textsc{ShapleyX}: A Python package for global sensitivity analysis with RS-HDMR.
\newblock \textit{Journal of Open Source Software} (submitted), 2026.

\bibitem{ShapleyXRepo}
F.~Bennett.
\newblock \textsc{ShapleyX} repository.
\newblock \url{https://github.com/frbennett/shapleyx}, v0.5.1, 2026.

\end{thebibliography}

\end{document}
